\section{Limitations and Extensions}

The proposed design provides a way to identify how experimentally varied migrant characteristics affect respondents' beliefs about labor-market and social integration outcomes and to examine how these beliefs are associated with respondents' personal experiences and information exposure. At the same time, several limitations are important for interpreting the proposed analyses and suggest directions for future research.\bigskip


\noindent\textbf{Endogeneity of experience and information measures.} A central limitation concerns the interpretation of the analyses in Section 4.2. While the vignette characteristics are randomly assigned, respondents' contact experiences and media exposure are not. Individuals may select into different forms of contact based on where they live, work, and socialize, and these factors may also be related to their beliefs about migrants. Similarly, media exposure may be selective, as individuals can choose migration-related information partly in accordance with their existing dispositions \citep{lueders2019refugees}. Reverse causality is also possible: respondents' prior beliefs may influence which interactions they seek, which information they consume, and how they evaluate these experiences and information. The inclusion of respondent-level controls can account for some observable differences between respondents but cannot eliminate these sources of endogeneity. The associations examined in H4--H6 should therefore not be interpreted as causal effects. Future research could address this limitation using longitudinal data that measure beliefs and experiences at different points in time or experimental designs that randomly vary the information respondents receive. Where suitable exogenous variation in opportunities for intergroup contact is available, this could also provide stronger evidence on the causal role of personal contact.\bigskip

\noindent\textbf{Measurement of memorable experiences and information.} A related limitation concerns H6a and H6b. The valence of respondents' most memorable personal experience or media information is reported retrospectively and therefore captures the respondent's current subjective evaluation rather than an objective measure of the original experience or information. What respondents remember and how they evaluate it may itself depend on their current beliefs, emotions, or the time that has passed since the experience occurred. Consequently, an association between memorable-experience valence and current beliefs does not establish that the remembered experience caused these beliefs. Moreover, these analyses are restricted to respondents for whom a memorable experience or piece of information comes to mind, and this group may systematically differ from respondents for whom nothing memorable is reported. A more direct extension would expose respondents to controlled information or experiences and measure both recall and beliefs after a delay. Such a design would make it possible to examine more directly whether differences in memory contribute to subsequent differences in beliefs.\bigskip

\noindent\textbf{Changes in beliefs over time and spillovers between migrant groups.} The cross-sectional design also cannot establish whether beliefs about Ukrainian or Syrian migrants have changed over time. This is particularly relevant because the large-scale arrival of Ukrainian refugees after the Russian invasion may have affected how other refugee groups are perceived. Current beliefs about Syrian refugees could be descriptively compared with findings from studies conducted before the arrival of Ukrainian refugees, including the period following the large arrival of Syrian refugees beginning in 2015. Such comparisons could indicate whether current patterns differ from those documented previously, but they would not show that the arrival of Ukrainian refugees caused these differences, since studies may differ in their samples, measures, and historical contexts. The possibility of spillovers could therefore be examined more directly in future research using comparable repeated surveys or panel data covering periods before and after the arrival of Ukrainian refugees. This could also build on existing research examining European attitudes toward refugees following the Russian invasion of Ukraine, such as \citet{moise2024european}.\bigskip

\noindent\textbf{Residence duration and integration over time.} The proposed vignettes do not vary migrants' year of arrival or duration of residence in Germany. As a result, the experiment cannot directly test whether otherwise comparable migrants are expected to become more integrated the longer they live in Germany. This also limits the connection between the vignette evaluations and the general beliefs measured by \(TimeLM_i\) and \(TimeSoc_i\). Including arrival year in the factorial design, however, would create difficulties because not all combinations of country of origin, refugee status, and arrival year are historically plausible. A future study could instead vary the stated duration of residence while restricting profiles to plausible combinations. This would allow beliefs about integration over time to be examined more directly and could show whether the perceived role of residence duration differs between migrant groups.\bigskip

\noindent\textbf{Scope and external validity of the vignette design.} The factorial vignette design necessarily simplifies migrant characteristics in order to isolate the effects of the experimentally varied attributes. In particular, the estimated country-of-origin effect captures the effect of presenting an otherwise comparable migrant as Ukrainian rather than Syrian, but it cannot identify which perceptions associated with country of origin account for this difference. Respondents may associate the two origins with differences in culture, religion, perceived similarity, migration history, or other characteristics that are not separately varied in the experiment. Similarly, stated beliefs about hypothetical profiles do not necessarily correspond to respondents' behavior or to migrants' actual integration outcomes. Future research could therefore vary additional theoretically relevant characteristics or compare stated beliefs with observed outcomes or behavioral measures.\bigskip

The focus on Ukrainian and Syrian migrants and on refugees and economic migrants also limits the extent to which the findings can be generalized to other migrant groups and migration pathways in Germany. Future studies could extend the design to additional countries of origin and reasons for migration, such as migration for education or family reasons. The set of outcome measures could likewise be expanded to other dimensions of economic and social integration. Such extensions would make it possible to examine whether the patterns identified in the present design are specific to the Ukrainian--Syrian comparison or generalize more broadly.

\noindent\textbf{Survey and measurement limitations.} Finally, several more general limitations arise from the proposed survey design. Contact frequency, experience quality, and media exposure are self-reported and may therefore be affected by recall and reporting error. Combining workplace and non-work contact into an overall contact measure may also conceal differences between forms of contact, although the two measures can additionally be examined separately. Moreover, the use of an online panel with demographic quotas can approximate the target population along the selected characteristics but does not ensure that respondents are representative of the German population in all respects. The migration context is also time- and country-specific, which limits generalization beyond contemporary Germany. Replication using other samples, migrant groups, countries, and time periods would therefore provide useful evidence on the broader generalizability of the findings.